%------------------------
% 이력서 — 오준석 (한국어)
% Based on tex_template/cv_template.tex formatting
% Source: docs/cv_ko.docx
% Compile: xelatex cv_ko.tex
%------------------------

\documentclass[a4paper,9pt]{extarticle}

\usepackage{kotex}
\usepackage{geometry}
\geometry{a4paper, margin=1in}
\usepackage{titlesec}
\usepackage{enumitem}
\usepackage{hyperref}
\usepackage{changepage}

\setlist{noitemsep, topsep=2pt, parsep=0pt, partopsep=0pt}
\titleformat{\section}{\large\bfseries}{}{0em}{}[\titlerule]
\titlespacing*{\section}{0pt}{\baselineskip}{\baselineskip}
\titleformat{\subsection}{\normalsize\bfseries}{}{0em}{}
\titlespacing*{\subsection}{0pt}{0.6\baselineskip}{0.3\baselineskip}

\newcommand{\researchentry}[4]{%
  \noindent\textbf{#1} \hfill \textit{#2}\\
  \textit{#3}
  \begin{itemize}
    #4
  \end{itemize}
  \vspace{0.4em}
}

\newcommand{\eduentry}[4]{%
  \noindent\textbf{#1} \hfill #4\\
  #2\\
  #3
  \vspace{0.3em}
}

\begin{document}
\pagestyle{empty}

\begin{center}
  \textbf{\Large Junseok Oh (오준석)}\\[4pt]
  \textbf{이력서 (Curriculum Vitae)}\\[4pt]
  (+82) 010-6640-0875 \,|\, \href{mailto:ohjs@sogang.ac.kr}{ohjs@sogang.ac.kr} \,|\, \href{https://june-oh.github.io/june-oh/}{june-oh.github.io/june-oh}
\end{center}

\section*{학력}
\eduentry{서강대학교}{컴퓨터공학과 박사 (지도교수: 김지환, 2026.08 졸업 예정)}{서울}{2022.03 -- 2026.08}
\begin{itemize}
  \item 학위논문: ``Adapter-Only Bridging of Frozen Speech Encoder and Frozen LLM for ASR''
\end{itemize}
\eduentry{서강대학교}{컴퓨터공학과 공학석사 (지도교수: 김지환)}{서울}{2017.09 -- 2019.08}
\begin{itemize}
  \item 학위논문: ``Korean Real-Time Automatic Transcription System Using Weakly Labeled Corpus''
\end{itemize}
\eduentry{서강대학교}{컴퓨터공학과 공학사}{}{2010.03 -- 2017.08}
\eduentry{의정부 고등학교}{고등학교 졸업}{경기도 의정부}{2007.03 -- 2010.02}

\section*{연구 경력}
\researchentry{대화 기반 설명가능성을 멀티모달로 제공하는 인공지능 기술 개발}{2022.04 -- 현재}{주관: 과학기술정보통신부(MSIT) / IITP}{
  \item 마비말장애(구음장애) 중증도 분류를 위한 AI 프레임워크 개발 및 멀티모달 설명 제공으로 진단 의사결정 지원
  \item 음향/언어 특성을 분석하는 설명가능 진단 모듈 구현 (중증도 수준에 대한 해석 가능한 피드백 제공)
  \item 복잡한 AI 모델 출력과 사용자 이해 가능한 설명 간 해석 격차를 대화형 멀티모달 환경에서 완화하는 방법 연구
}

\researchentry{End-to-End 한국어 음성인식 프로젝트}{2024 -- 2025.04}{산학협력: 스마일게이트}{
  \item Hybrid FastConformer RNNT + CTC 기반의 범용 한국어 ASR 시스템 개발
  \item 인식 정확도 향상 및 지연 감소를 위한 cache-aware streaming 기법 구현
  \item 게임 도메인 용어 등 도메인 특화 어휘에 대한 컨텍스트 바이어싱(context biasing) 적용으로 인식 성능/강건성 향상
}

\researchentry{전화망(8kHz) 환경 End-to-End 음성인식}{2024.04 -- 2024.12}{산학협력: 롯데이노베이트}{
  \item 8kHz 전화망 데이터에 최적화된 스트리밍/비스트리밍 한국어 ASR 파이프라인 개발 (FastConformer-CTC)
  \item 도메인 시프트에 대응하기 위한 동적 컨텍스트 바이어싱 모듈 구현 (단어 수준 정확도 개선)
}

\researchentry{(외국인을 위한) 2024년 한국어 평가 말하기, 쓰기 자동 채점 연구}{2024.05 -- 2024.12}{주관: 문화체육관광부}{
  \item Wav2Vec 기반 멀티 태스크 학습으로 발음(pronunciation), 유창성(fluency), 내용(content)을 공동 모델링하는 L2-한국어 말하기 평가 프레임워크 개발
  \item Conformer-CTC ASR 출력을 대규모 언어모델(LLaMa)과 결합해 다측면 자동 채점/점수 예측
}

\researchentry{무분별적 유해성 미디어 배포 문제 해결을 위한 지능적 분석/분류 기반 콘텐츠 등급분류 기술 개발}{2022 -- 2024}{주관: 과학기술정보통신부(MSIT)}{
  \item 자동 영상 등급 판정 프레임워크 내 음성인식 및 오디오 분석, 주요 멀티미디어 이벤트를 위한 음향 이벤트 검출(sound-event detection) 모델 설계
  \item 도메인 영상 코퍼스에 맞춰 OpenAI Whisper ASR을 파인튜닝하여 다양한 음향 환경에서 전사 강건성 향상
}

\researchentry{(외국인을 위한) 2023년 한국어 평가 말하기, 쓰기 자동 채점 연구}{2023.05 -- 2023.12}{주관: 문화체육관광부}{
  \item Conformer-CTC ASR 출력 + BERT 기반 의미(semantic) 점수화를 결합해 L2-한국어 평가 파이프라인 구축
  \item ASR 전사로부터 발음 정확도, 발화 속도(speech rate), 구문적 정확성(syntactic correctness)을 정량화하는 알고리즘 개발
}

\researchentry{비디오 튜링 테스트를 통과할 수준의 비디오 스토리 이해 기반의 질의응답 기술 개발}{2017.09 -- 2019.12}{주관: 과학기술정보통신부(MSIT)}{
  \item Kaldi 문장 단위(sentence-level) 디코더를 수정하여 실시간 비디오 QA에서 RT(real-time factor) 1.0 미만 달성
  \item 타깃 비디오 시리즈에서 도메인 특화 오디오/텍스트 코퍼스를 수집·정제하고 음향/언어 모델 최적화
  \item 신규 코퍼스를 활용한 음향모델 재학습 및 언어모델 업데이트로 복잡 내러티브에서 QA 정확도 개선
}

\section*{논문}
\subsection*{International Journal}
\begin{itemize}
  \item J. Oh, J. Nam, and J.-H. Kim, ``HiTCA: Fusing Hierarchical Text and Contextual Audio for Accurate VCR,'' \textit{EURASIP Journal on Audio, Speech, and Music Processing}, under review. [SCIE]
  \item S. Ma, J. Oh, M. Kim, and J.-H. Kim, ``Survey on Deep Learning-based Speech Technologies in Voice Chatbot Systems,'' \textit{KSII Transactions on Internet \& Information Systems (TIIS)}, vol.~19, no.~5, pp.~1406--1440, 2025. [SCIE]
  \item J. Oh, E. Cho, and J.-H. Kim, ``Integration of WFST language model in pre-trained Korean E2E ASR model,'' \textit{KSII Transactions on Internet and Information Systems (TIIS)}, vol.~18, no.~6, pp.~1692--1705, 2024. [SCIE]
  \item S. Seo, J. Oh, E. Cho, H. Park, G. Kim, and J.-H. Kim, ``TP-MobNet: A Two-pass Mobile Network for Low-complexity Classification of Acoustic Scene,'' \textit{Computers, Materials \& Continua}, vol.~73, no.~2, 2022. [SCIE]
  \item M. Lim, D. Lee, H. Park, Y. Kang, J. Oh, J.-S. Park, G.-J. Jang, and J.-H. Kim, ``Convolutional neural network based audio event classification,'' \textit{KSII Transactions on Internet and Information Systems (TIIS)}, vol.~12, no.~6, pp.~2748--2760, Jun. 2018. [SCIE]
\end{itemize}

\subsection*{International Conference}
\begin{itemize}
  \item J. Oh and J.-H. Kim, ``Adapter-Only Bridging of Frozen Speech Encoder and Frozen LLM for ASR,'' in \textit{Advances in Knowledge Discovery and Data Mining (PAKDD 2026)}, Lecture Notes in Computer Science, vol.~16600, pp.~251--262, Springer, Singapore, 2026. [Oral]
  \item J. Oh and J.-H. Kim, ``SEAM: Bridging the Temporal-Semantic Granularity Gap for LLM-based Speech Recognition,'' in \textit{Findings of the Association for Computational Linguistics: EACL 2026}, pp.~2135--2144, 2026.
  \item J. Oh, H. Park, and J.-H. Kim, ``Speech Intelligibility Prediction of Dysarthria Using Deep Convolutional Networks,'' in \textit{Proc.\ APIC-IST 2023}, pp.~236--237, 2023.
  \item M. Kim, J. Oh, and J.-H. Kim, ``Automated Dysarthria Severity Classification Using Diadochokinetic test and Speech Intelligibility Based on LightGBM,'' in \textit{Proc.\ APIC-IST 2023}, pp.~12--13, 2023.
  \item S. Seo, M. Lim, D. Lee, H. Park, J. Oh, D. J. Rim, and J.-H. Kim, ``Environmental noise robustness for Korean fricatives using speech enhancement generative adversarial networks,'' in \textit{Proc.\ IEEE Int.\ Conf.\ Big Data and Smart Computing (BigComp)}, pp.~1--4, 2019.
  \item S. Seo, D. J. Rim, M. Lim, D. Lee, H. Park, J. Oh, C. Kim, and J.-H. Kim, ``Shortcut connections based deep speaker embeddings for end-to-end speaker verification system,'' in \textit{Proc.\ Interspeech}, pp.~2928--2932, 2019.
\end{itemize}

\subsection*{Domestic Journal}
\begin{itemize}
  \item 이정필, 장재후, 김지현, 김민섭, 김성준, 김민서, 김하영, 오준석, 정원, 김장연 외, ``음성에 기반한 마비말장애 진단과 설명이 가능한 시스템,'' \textit{정보과학회지}, 42(4): 45--56, 2024. [KSCI/KCI]
  \item H. Park, Y. Kang, M. Lim, D. Lee, J. Oh, and J.-H. Kim, ``LFMMI-based acoustic modeling by using external knowledge,'' \textit{The Journal of the Acoustical Society of Korea}, vol.~38, no.~5, pp.~607--613, 2019. [KSCI/KCI]
\end{itemize}

\section*{수상}
\begin{itemize}
  \item \textbf{2023년 한국어 AI 경진대회} (한국지능정보사회진흥원, NIA) --- Track2-1, 상담 음성인식: 오준석, 김민서, 남주형 (team ``상담ONE''), 장려상
  \item \textbf{2022 한국어 인공지능 경진대회} (한국지능정보사회진흥원, NIA) --- 기업현안(회의음성): 오준석, 김하영 (team ``SGCSE''), 최우수상/네이버 대표 (1위)
  \item \textbf{2021 음절인식률 측정 알고리즘 개발 대회} (KT alpha) --- 박호성, 오준석, 조은수 (team ``검은사케동''), 최우수상 (1위)
\end{itemize}

\section*{특허}
\begin{itemize}
  \item KR 10-2699607 (B1) --- \textit{Corpus Construction Service Provision Server and Method} (코퍼스 구축 서비스 제공 서버 및 방법).\\
  발명자: Ji-Hwan Kim; Junseok Oh; Hosung Park.\\
  출원인: 서강대학교 산학협력단 (Sogang University Industry--Academic Cooperation Foundation).\\
  출원: 2020.04.29 (KR 10-2020-0052570); 공개: 2021.11.08 (KR 10-2021-0133667); 등록: 2024.08.22; 공보: 2024.08.26.
\end{itemize}

\section*{자격증}
\begin{itemize}
  \item NVIDIA Deep Learning Institute Certificate, \textit{Building Conversational AI Applications}, 2022
\end{itemize}

\section*{기술 역량}
\begin{itemize}
  \item \textbf{Programming Languages \& Scripting:} Python
  \item \textbf{Deep Learning Frameworks:} PyTorch, NVIDIA NeMo, Hugging Face Transformers, PEFT
  \item \textbf{E2E ASR Toolkits:} Kaldi
  \item \textbf{Other Tools \& Libraries:} KenLM
\end{itemize}

\end{document}
